\section{La brecha entre el desempeño en eval y la capacidad real}

\subsection{La esquizofrenia del modelo}

Ilya señala un fenómeno desconcertante en los modelos actuales: tienen un desempeño excelente en eval, pero en el uso real siguen cometiendo errores elementales. Lo ilustra con un ejemplo de vibe coding: el modelo corrige un bug, pero introduce un segundo bug; le señalas el segundo bug, y al corregirlo vuelve a introducir el primero, pudiendo quedar atrapado en un ciclo infinito entre ambos bugs.

\begin{importantbox}{Dos explicaciones posibles}
\begin{enumerate}
    \item \textbf{El RL hace que el modelo se enfoque de manera excesivamente unidireccional}: el entrenamiento con RL puede hacer que el modelo se vuelva "too single-minded" en ciertos aspectos; aunque también refuerza otras capacidades, el juicio básico puede verse afectado
    \item \textbf{El diseño del entorno de RL se inspira en el eval}: al diseñar los entornos de entrenamiento con RL, los equipos de investigación toman inspiración del eval sin darse cuenta ("quiero que el modelo tenga buen desempeño en este eval, ¿qué tipo de entrenamiento con RL puede ayudar aquí?"), lo que en realidad hace que el modelo termine "teaching to the test"
\end{enumerate}
\end{importantbox}

\subsection{La analogía de las 10,000 horas contra las 100 horas}

\begin{importantbox}{Competidor entrenado contra competidor talentoso}
Supongamos que dos estudiantes estudian programación competitiva:
\begin{itemize}
    \item El estudiante A practicó 10,000 horas, resolvió incontables problemas y memorizó todas las técnicas de demostración: se convirtió en un competidor de primer nivel
    \item El estudiante B solo practicó 100 horas, pero tuvo un desempeño igualmente sobresaliente
\end{itemize}
¿Cuál de los dos tiene mejores perspectivas de desarrollo profesional? \textbf{El segundo}. Porque la capacidad del estudiante B proviene de una capacidad de generalización más profunda, no de un entrenamiento excesivo en un dominio específico. Los modelos de IA actuales se parecen más al estudiante A, e incluso de forma más extrema, porque no solo usamos todos los problemas de competencias, sino que además hicimos aumento de datos para generar más problemas.
\end{importantbox}

\begin{warningbox}{El verdadero reward hacking}
Ilya sugiere: el verdadero reward hacking no lo hace el modelo, sino los \textbf{investigadores humanos}. Al prestar demasiada atención al eval, alinean sin darse cuenta el entorno de entrenamiento con RL con las métricas del eval, lo que produce un sesgo sistemático entre el desempeño en eval y la capacidad real.
\end{warningbox}

\subsection{Resumen de la sección}

Existe una brecha considerable entre el desempeño en eval y la capacidad en el mundo real. Las causas posibles son que el entrenamiento con RL se enfoca en exceso en dominios verificables, y que los investigadores humanos, sin darse cuenta, hacen "teaching to the test". Los modelos se parecen más a competidores sobreentrenados que a verdaderos generalistas.
